%=============================================================================
% Appendices
%=============================================================================
\appendix

\section{Capability Execution Sequence Diagram}
\label{appendix:sequence}

\subsection{Complete Execution Flow}

The following sequence diagram illustrates the complete flow from agent intent to capability execution:

\begin{figure}[htbp]
\centering
\begin{tikzpicture}[
    participant/.style={
        draw,
        rectangle,
        minimum width=1.5cm,
        minimum height=0.55cm,
        align=center,
        font=\footnotesize
    },
    arrow/.style={
        -{Stealth[length=2mm, width=1.5mm]},
        thick
    },
    returnarrow/.style={
        -{Stealth[length=2mm, width=1.5mm]},
        thick,
        dashed
    },
    label/.style={font=\scriptsize}
]

% Participants
\node[participant, fill=blue!12] (agent) at (0, 0) {Agent\\Layer};
\node[participant, fill=green!12] (orch) at (2.8, 0) {Orchest.\\Layer};
\node[participant, fill=yellow!15] (contact) at (5.6, 0) {Contact\\Layer};
\node[participant, fill=orange!18] (exec) at (8.4, 0) {Execution\\Runtime};
\node[participant, fill=purple!12] (prov) at (11.2, 0) {Provider\\API};

% Vertical lines
\foreach \x in {0, 2.8, 5.6, 8.4, 11.2} {
    \draw[thick, gray!50] (\x, -0.45) -- (\x, -7.5);
}

% Messages (numbered)
\draw[arrow] (0, -1) -- node[above, label] {1. Task Intent} (2.8, -1);
\draw[arrow] (2.8, -1.6) -- node[above, label] {2. Capability Req.} (5.6, -1.6);
\draw[arrow] (5.6, -2.2) -- node[above, label] {3. Query Registry} (5.6, -2.5);
\draw[arrow] (5.6, -3.1) -- node[above, label] {4. Available Providers} (5.6, -3.4);
\draw[arrow] (5.6, -4.0) -- node[above, label] {5. Routing Decision} (5.6, -4.3);
\draw[arrow] (5.6, -4.9) -- node[above, label] {6. Exec Directive} (8.4, -4.9);
\draw[arrow] (8.4, -5.5) -- node[above, label] {7. API Call} (11.2, -5.5);
\draw[returnarrow] (11.2, -6.0) -- node[above, label] {8. Response} (8.4, -6.0);
\draw[returnarrow] (8.4, -6.5) -- node[above, label] {9. Normalized} (5.6, -6.5);
\draw[returnarrow] (5.6, -7.0) -- node[above, label] {10. Capability} (2.8, -7.0);
\draw[returnarrow] (2.8, -7.5) -- node[above, label] {11. Task Result} (0, -7.5);

\end{tikzpicture}
\caption{Complete execution flow from agent intent to capability execution.}
\label{fig:sequence-complete}
\end{figure}

\subsection{Error Handling Flow}

When execution fails, the error propagation path:

\begin{figure}[htbp]
\centering
\begin{tikzpicture}[
    participant/.style={
        draw,
        rectangle,
        minimum width=1.8cm,
        minimum height=0.55cm,
        align=center,
        font=\footnotesize
    },
    arrow/.style={
        -{Stealth[length=2mm, width=1.5mm]},
        thick
    },
    returnarrow/.style={
        -{Stealth[length=2mm, width=1.5mm]},
        thick,
        dashed,
        red!70!black
    },
    label/.style={font=\scriptsize}
]

% Participants
\node[participant, fill=green!12] (orch) at (0, 0) {Orchest.\\Layer};
\node[participant, fill=yellow!15] (contact) at (4, 0) {Contact\\Layer};
\node[participant, fill=orange!18] (exec) at (8, 0) {Execution\\Runtime};
\node[participant, fill=purple!12] (prov) at (12, 0) {Provider\\API};

% Vertical lines
\foreach \x in {0, 4, 8, 12} {
    \draw[thick, gray!50] (\x, -0.45) -- (\x, -6.5);
}

% Messages
\draw[arrow] (0, -1) -- node[above, label] {Execution Directive} (8, -1);
\draw[arrow] (8, -1.6) -- node[above, label] {API Call} (12, -1.6);
\draw[arrow] (8, -2.2) -- node[above, label] {Request} (12, -2.2);
\draw[returnarrow] (12, -3.0) -- node[above, label, text=red!70!black] {Error (429, 502)} (8, -3.0);

% Decision box
\node[draw, rectangle, rounded corners, fill=orange!10, font=\scriptsize, minimum width=2.5cm] at (8, -4) {Micro-Retry?};
\node[font=\scriptsize, text=gray] at (8, -4.5) {(if enabled \& transient)};

\draw[returnarrow] (8, -5.2) -- node[above, label, text=red!70!black] {Execution Error} (4, -5.2);
\node[font=\scriptsize, text=gray] at (2, -5.7) {Policy Decision: retry/fallback/abort};

\end{tikzpicture}
\caption{Error handling flow with micro-retry decision and policy escalation.}
\label{fig:sequence-error}
\end{figure}

%-----------------------------------------------------------------------------
\section{Capability Definition Examples}
\label{appendix:examples}

\subsection{Text Generation Capability}

\begin{lstlisting}[style=yamlstyle]
name: generate_text
version: "1.0"
description: Generate text based on a prompt
provider: openai
model: gpt-5.4
input_schema:
  type: object
  properties:
    prompt:
      type: string
    max_tokens:
      type: integer
      default: 1000
    temperature:
      type: number
      default: 0.7
output_schema:
  type: object
  properties:
    text:
      type: string
    tokens_used:
      type: integer
execution:
  endpoint: /v1/chat/completions
  method: POST
  timeout_ms: 30000
\end{lstlisting}

\subsection{Web Search Capability}

\begin{lstlisting}[style=yamlstyle]
name: search_web
version: "1.0"
description: Search the web for information
provider: search_service
input_schema:
  type: object
  properties:
    query:
      type: string
    num_results:
      type: integer
      default: 10
output_schema:
  type: object
  properties:
    results:
      type: array
      items:
        type: object
        properties:
          url:
            type: string
          title:
            type: string
          snippet:
            type: string
execution:
  endpoint: /v1/search
  method: GET
\end{lstlisting}

%-----------------------------------------------------------------------------
\section{Runtime Pseudocode}
\label{appendix:pseudocode}

\subsection{Execute Capability}

\begin{lstlisting}[style=pythonstyle]
def execute_capability(request: CapabilityRequest) -> Result:
    """
    Execute a capability request.

    Args:
        request: Standardized capability request

    Returns:
        Normalized execution result
    """
    # 1. Validate request
    validate_request(request)

    # 2. Translate to provider-specific format
    provider_request = translate_request(request)

    # 3. Invoke provider API
    try:
        provider_response = invoke_provider(
            endpoint=request.provider.endpoint,
            data=provider_request
        )
    except ConnectionError as e:
        # Propagate error (no retry logic)
        raise ExecutionError(
            capability=request.capability,
            error_type="connection",
            message=str(e)
        )

    # 4. Normalize response
    result = normalize_response(
        capability=request.capability,
        provider_response=provider_response
    )

    # 5. Return standardized result
    return result
\end{lstlisting}

\subsection{Request Translation}

\begin{lstlisting}[style=pythonstyle]
def translate_request(request: CapabilityRequest) -> dict:
    """
    Translate standardized request to provider-specific format.

    Args:
        request: Standardized capability request

    Returns:
        Provider-specific request payload
    """
    capability_manifest = load_manifest(request.capability)
    provider_adapter = get_adapter(request.provider.name)

    return provider_adapter.translate(
        capability=request.capability,
        input_data=request.input,
        manifest=capability_manifest
    )
\end{lstlisting}
